\chapter{Nếu phải giữ lại một ý}
\begin{tinhhuongbox}{Tình huống | Situation}
\songngu{Giữa tiết, vài ánh mắt bắt đầu ra cửa và vài cái đầu bắt đầu cúi xuống.}{In the middle of the lesson, a few eyes start moving toward the door and a few heads start dropping.}
\songngu{Đó là lúc phải kéo suy nghĩ quay lại trước khi kéo kỷ luật quay lại.}{That is when you should pull thinking back before pulling discipline back.}
\end{tinhhuongbox}
\begin{noitrenlop}
\songngu{Nếu chỉ được giữ lại một ý trong ba phút vừa rồi, em giữ ý nào?}{If you could keep only one idea from the last three minutes, which one would you keep?}
\end{noitrenlop}
\begin{vihieuqua}
\songngu{Câu hỏi này cắt mạch trôi mà không cắt mặt lớp.}{This question cuts the drifting without cutting the class's face.}
\end{vihieuqua}
\begin{englishpocket}
\songngu{Cụm nên nhớ: \textit{Which one would you keep?}}{Phrase to keep: \textit{Which one would you keep?}}
\end{englishpocket}
\chapter{Lớp mình đang mệt hay đang chán?}
\begin{tinhhuongbox}{Tình huống | Situation}
\songngu{Nhiều tiết trôi không phải vì nội dung dở, mà vì năng lượng lớp xuống không được gọi tên.}{Many lessons drift not because the content is bad, but because the class's low energy is never named.}
\songngu{Gọi đúng vấn đề sẽ cứu nhịp nhanh hơn đoán mò.}{Naming the problem correctly saves the rhythm faster than guessing.}
\end{tinhhuongbox}
\begin{noitrenlop}
\songngu{Lớp mình đang mệt hay đang chán, để tôi cứu cho đúng bệnh?}{Are we tired or bored, so I can treat the right problem?}
\end{noitrenlop}
\begin{vihieuqua}
\songngu{Lớp thường thích câu này vì nó vừa thật vừa có chút hài.}{Students usually like this line because it feels honest and slightly funny.}
\end{vihieuqua}
\begin{englishpocket}
\songngu{Cụm nên nhớ: \textit{Treat the right problem.}}{Phrase to keep: \textit{Treat the right problem.}}
\end{englishpocket}
\chapter{Đứng lên chọn phe}
\begin{tinhhuongbox}{Tình huống | Situation}
\songngu{Có lúc cứu lớp không phải bằng thêm lời mà bằng đổi tư thế.}{Sometimes you save a class not by adding more words, but by changing posture.}
\songngu{Một chuyển động ngắn có thể thay cả nhịp tiết.}{One short movement can change the rhythm of the whole lesson.}
\end{tinhhuongbox}
\begin{noitrenlop}
\songngu{Ai đồng ý đứng bên trái, ai còn phân vân đứng bên phải, đứng lên cho não đỡ ngủ.}{If you agree, stand on the left. If you are unsure, stand on the right. Stand up so your brain stops sleeping.}
\end{noitrenlop}
\begin{vihieuqua}
\songngu{Cơ thể nhúc nhích một chút thì đầu óc cũng nhúc nhích theo.}{When the body moves a little, the mind tends to move with it.}
\end{vihieuqua}
\begin{englishpocket}
\songngu{Cụm nên nhớ: \textit{If you are unsure.}}{Phrase to keep: \textit{If you are unsure.}}
\end{englishpocket}
\chapter{Đổi tốc độ chứ chưa cần đổi bài}
\begin{tinhhuongbox}{Tình huống | Situation}
\songngu{Có tiết không cần cứu bằng nội dung mới, chỉ cần cứu bằng nhịp mới.}{Some lessons do not need new content; they only need a new rhythm.}
\songngu{Tăng tốc rất ngắn rồi hạ lại là một cách kéo lớp tỉnh nhanh.}{A very short speed-up followed by a calm return is one way to wake the class quickly.}
\end{tinhhuongbox}
\begin{noitrenlop}
\songngu{Ba mươi giây tốc độ cao: mỗi bàn nói đúng một câu, không vòng vo như họp tổ dân phố.}{Thirty seconds at high speed: each desk says exactly one sentence, no winding speeches like a neighborhood meeting.}
\end{noitrenlop}
\begin{vihieuqua}
\songngu{Lớp bật cười vì ví von, rồi tự rút gọn lời nói lại.}{The class laughs at the comparison and then naturally shortens their answers.}
\end{vihieuqua}
\begin{englishpocket}
\songngu{Cụm nên nhớ: \textit{Exactly one sentence.}}{Phrase to keep: \textit{Exactly one sentence.}}
\end{englishpocket}
\chapter{Đừng để năm bạn đầu tiên kéo theo cả lớp}
\begin{tinhhuongbox}{Tình huống | Situation}
\songngu{Khi năm bạn đầu tiên bắt đầu lơ là, phần còn lại thường đi theo rất nhanh.}{When the first five students start drifting away, the rest usually follow fast.}
\songngu{Phải chặn ngay ở điểm đầu lan rộng.}{You have to stop it right at the point where it starts spreading.}
\end{tinhhuongbox}
\begin{noitrenlop}
\songngu{Năm bạn đầu tiên đang chuẩn bị kéo cả lớp đi xa, kéo lại giúp tôi một câu nào.}{The first five students are about to drag the whole class away, bring us back with one sentence.}
\end{noitrenlop}
\begin{vihieuqua}
\songngu{Câu này không gọi tên cá nhân nhưng đủ để cả lớp biết tín hiệu đã được thấy.}{This line does not name individuals, but it tells the class clearly that the signal has been noticed.}
\end{vihieuqua}
\begin{englishpocket}
\songngu{Cụm nên nhớ: \textit{Bring us back.}}{Phrase to keep: \textit{Bring us back.}}
\end{englishpocket}